% 第8章第4节：本章小结（同时收束“函数空间”部分）。

\section{本章小结}
\label{sec:apx-summary}

章首提出的三个问题，分别对应逼近对象的三个层次。

对一个给定的函数，逼近的含义取决于固定了什么：定义域、误差范数和生成单元。Bernstein多项式、线性样条、固定宽度的高斯核截面和单隐藏层网络都能一致逼近紧集上的任意连续目标，但这些都是可变预算下的存在性结论，次数、结点、中心或单元数随精度增加。误差随预算下降的速度另需光滑性：连续模给出$\omega_f(h)$量级，二阶导数有界给出$h^2$量级，折线$f_*$只有一阶。稠密性还不等于能稳定地算出系数：高斯核和的系数可能随精度发散，在固定核范数预算下甚至无法任意逼近$f_*$；网络参数也不可识别。

对多变量函数，乘积结构提供了组合的方式。在可分的$L^2$空间中，单变量正交基的乘积构成乘积测度下的正交基，张量积网格却使系数个数按$m^d$增长。低秩表示只保留少数乘积项；在正交基下，最佳$r$项$L^2$近似由系数矩阵的截断奇异值分解给出，非正交基须先计入Gram矩阵。项数、因子幅度与误差是不同的量，三组以上变量时最佳低秩近似甚至可能不存在。

对整个函数类，量词从“对每个目标存在近似”变为“存在一套代表适用于全体”。一致距离下，紧度量空间上的函数族相对紧，当且仅当一致有界且等度连续；核范数有界球因点值界与核截面差界而在$C(X)$中紧，却不在核范数下紧。覆盖数计量所需代表的个数，它依赖所选的距离：经验距离下的有限覆盖不能推出全域覆盖。常值族、有限维参数族与Lipschitz类的估计表明，覆盖数随幅度、参数维数、正则性和定义域维数变化，而不由参数个数单独决定。

至此，本部分建立了研究函数的几套工具：第~\ref{chap:function-spaces}章的范数、完备性与正交几何，第~\ref{chap:linear-functionals-operators}章与第~\ref{chap:operator-spectra-equations}章的泛函、算子、谱与方程，第~\ref{chap:weak-derivatives-sobolev}章以弱导数刻画的正则性，第~\ref{chap:rkhs}章以点值观测构造的核空间，以及本章的有限逼近与覆盖。它们回答的是候选函数在什么空间中、怎样度量、怎样表示和能逼近到什么程度。后续的优化问题在这些候选空间上加入目标函数与约束，研究最优解的存在、最优条件和求解方法，可参考\BookChapterRef{chap:optimization}{卷三“优化问题的定义与基本性质”章}；信号处理加入时间、频率和采样机制；统计与学习理论加入随机数据和误差保证。它们沿用这里已经确定的范数、表示与覆盖结论，只在各自的问题中补充新的结构。
